\chapter{Quantum Linear Systems}
\label{ch:lec10}

\begin{objectives}
\begin{itemize}[leftmargin=1.4em,nosep]
  \item Formulate the quantum linear system problem (QLSP) and its parameters.
  \item Solve QLSP with HHL and with QSVT.
  \item Explain adiabatic and Zeno-based near-optimal solvers.
  \item Compare the complexity landscape across methods.
\end{itemize}
\end{objectives}

The quantum numerical linear algebra of the last four chapters culminates in the
oldest and most-studied ``killer application'' of the whole program: solving a
linear system $A\ket x\propto\ket b$ on a quantum computer. The input models of
\cref{ch:lec06}, the LCU arithmetic of \cref{ch:lec07}, the qubitization of
\cref{ch:lec08}, and above all the matrix-inversion polynomial $x\mapsto1/(\kappa x)$
of \cref{ch:lec09} were each introduced with this target in mind. What makes the
problem subtle is not implementing $A^{-1}$ on paper---that is a single QSVT sweep---but
controlling the cost as a function of two adversarial parameters: the target error
$\epsilon$ and, decisively, the \emph{condition number} $\kappa=\norm A\norm{A^{-1}}$,
which measures how nearly singular $A$ is and which no generic quantum solver can
escape (\cref{sec:lec10-1}). We tell the story historically and by increasing
sophistication. The Harrow--Hassidim--Lloyd (HHL) algorithm (\cref{sec:lec10-2}) makes
the mechanism transparent---phase-estimate the eigenvalues of $A$, rotate each by
$1/\lambda$, uncompute, postselect---and exposes the two sources of inefficiency
(precision-limited QPE and small success probability) that drive its
$\bigo(\kappa^2/\epsilon)$ cost. QSVT (\cref{sec:lec10-3}) replaces the eigenvalue
arithmetic by a direct degree-$\bigo(\kappa\log(1/\epsilon))$ polynomial, improving the
$\epsilon$ dependence to logarithmic. Both, even after amplitude amplification, retain
a quadratic $\kappa$ dependence in the worst case; the adiabatic and quantum Zeno
solvers (\cref{sec:lec10-4}) recast QLSP as tracing a zero-eigenvalue eigenpath from
$\ket b$ to $\ket x$ and reach the near-optimal and optimal $\bigo(\kappa\,\polylog)$
and $\bigo(\kappa\log(1/\epsilon))$ scalings. We close (\cref{sec:lec10-5}) with the
complexity landscape and with the practical fact that $\kappa$, not the exponential
Hilbert-space dimension, governs everything---motivating preconditioning, whose
quantum incarnation (block-encoding normalization, the $(A+B)$ pipeline, and
multilevel BPX preconditioners) is where genuine speedups over classical iterative
solvers are won or lost. Throughout, $A\in\C^{N\times N}$ with $N=2^n$ is accessed by a
block encoding $U_A\in\mathrm{BE}_{\alpha,m}(A)$ and $\ket b$ by a state-preparation
unitary $U_b$ with $U_b\ket{0^n}=\ket b$.

\section{Problem setup: condition number and error}
\label{sec:lec10-1}

\paragraph{The problem and its output model.}
Classically, given $A\in\C^{N\times N}$ and $b\in\C^N$, one seeks $x=A^{-1}b$. The
\emph{quantum} linear system problem asks only for the \emph{direction} of the
solution, encoded as a normalized state.

\begin{definition}[Quantum linear system problem]
\label{def:lec10-qlsp}
Let $A\in\C^{N\times N}$ be invertible, accessed by a block encoding
$U_A\in\mathrm{BE}_{\alpha,m}(A)$, and let $\ket b$ be a normalized state prepared by
$U_b\ket{0^n}=\ket b$. Given a tolerance $\epsilon\in(0,1)$, the \emph{QLSP} is to
prepare a state $\ket{\tilde x}$ with
\begin{equation}
  \norm{\ket{\tilde x}-\ket x}\le\epsilon,\qquad
  \ket x=\frac{A^{-1}\ket b}{\norm{A^{-1}\ket b}}.
  \label{eq:lec10-qlsp}
\end{equation}
The \emph{query complexity} is the number of calls to $U_A,U_A\dg$ and to $U_b,U_b\dg$.
A solver may fail with probability bounded away from $1$, signaling failure through a
flag qubit.
\end{definition}

Two normalizations are standard. First, we always rescale so that $\norm A\le1$; after
the block-encoding subnormalization this places the singular values of $A$ in
$[1/\kappa,1]$, where
\begin{equation}
  \kappa=\kappa(A)\defeq\norm A\,\norm{A^{-1}}
        =\frac{\sigma_{\max}(A)}{\sigma_{\min}(A)}\ge1
  \label{eq:lec10-kappa}
\end{equation}
is the condition number. Second, without loss of generality we may take $A$
\emph{Hermitian positive definite}: a general $A$ is handled by the Hermitian dilation
\begin{equation}
  \widetilde A=\begin{pmatrix}0&A\\ A\dg&0\end{pmatrix},\qquad
  \widetilde A\begin{pmatrix}0\\ x\end{pmatrix}=\begin{pmatrix}b\\ 0\end{pmatrix},
  \label{eq:lec10-dilation}
\end{equation}
which has the same singular values (hence the same $\kappa$) and returns the solution
$x$ in the lower block (\cref{ch:lec09}, \eqref{eq:lec09-dilation}). We freely use both
viewpoints below.

\paragraph{Why $\kappa$ and $\epsilon$ drive the cost.}
The map $x\mapsto1/x$ must be resolved on the whole spectral interval $[1/\kappa,1]$.
Near the small singular value $\sigma_{\min}=1/\kappa$ the inverse is as large as
$\kappa$, so any device that outputs $A^{-1}$ at unit subnormalization can only encode
$A^{-1}/\kappa$: the solution appears with amplitude $\sim1/\kappa$ inside a
larger state, and \emph{postselecting} it succeeds with probability $\sim1/\kappa^2$
in the worst case. This is the origin of every $\kappa$ (or $\kappa^2$) in the
complexities below. The error $\epsilon$ enters because the inverse must be
approximated; how gently it enters (as $1/\epsilon$ or $\log(1/\epsilon)$) is exactly
what distinguishes the methods. That $\kappa$ cannot be removed is a theorem, not an
artifact of any particular algorithm.

\begin{proposition}[Condition-number lower bound]
\label{prop:lec10-lb}
Any quantum algorithm that solves the QLSP of \cref{def:lec10-qlsp} for all inputs with
a given $\kappa$ must make $\Omega(\kappa)$ queries to the oracle for $A$.
\end{proposition}

\begin{proof}[Proof idea]
One embeds a Hamiltonian-simulation instance into a linear system so that a
sublinear-in-$\kappa$ QLSP solver would simulate $\ee^{-\ii Ht}$ for time $t\sim\kappa$
with $o(\kappa)$ queries, violating the ``no fast-forwarding'' theorem of
\cref{ch:lec12} for generic Hamiltonians. Thus the linear scaling in $\kappa$ of the
optimal solvers (\cref{sec:lec10-4}) is best possible for generic $A$; special
structure (e.g.\ positive definiteness, exploited classically by conjugate gradient to
get $\sqrt\kappa$) can beat it only in restricted settings.
\end{proof}

\paragraph{Reading out the answer.}
The QLSP outputs the \emph{normalized} state $\ket x$; the norm $\norm{A^{-1}\ket b}$ is
recovered from the success probability of the flagged block encoding, since a QLSP
solver realizes $\widetilde U\in\mathrm{BE}_{\gamma,m}(A^{-1},\epsilon)$ and the
unflagged vector $\ket{\tilde x\!\succ}=\bra{0^m}\widetilde U\ket{0^m}\ket b$ satisfies
$\norm{x-\gamma\ket{\tilde x\!\succ}}\le\epsilon$, so $\gamma\norm{\ket{\tilde
x\!\succ}}\approx\norm x$. Downstream one typically wants an observable
$\bra x O\ket x$, estimated by a Hadamard or swap test on $\ket x$, or an overlap
$\braket{y}{x}$; these are the ``end-to-end'' tasks in which a QLSP solver is a
subroutine. Crucially, the quantum solver never writes out the $N$ entries of $x$; the
hoped-for exponential advantage lives entirely in the ability to prepare $\ket x$ and
sample a \emph{scalar} functional of it in time $\poly(n)$ rather than $\poly(N)$.

\section{Phase estimation and HHL}
\label{sec:lec10-2}

The Harrow--Hassidim--Lloyd algorithm was the first quantum QLSP solver. Take $A$
Hermitian with eigendecomposition $A\ket{v_j}=\lambda_j\ket{v_j}$, accessed through the
Hamiltonian-evolution oracle $U=\ee^{\ii2\pi A}$, and expand
$\ket b=\sum_j c_j\ket{v_j}$. Then
\begin{equation}
  A^{-1}\ket b=\sum_j\frac{c_j}{\lambda_j}\ket{v_j},
  \label{eq:lec10-hhl-target}
\end{equation}
so it suffices to multiply the $j$-th eigencomponent by $1/\lambda_j$. HHL does this in
four steps, sketched in the circuit of \cref{fig:lec10-hhl}.

\paragraph{The four steps.}
Assume for the analysis that each $\lambda_j\in(0,1)$ has an exact $d$-bit binary
expansion (this idealization isolates the mechanism; the finite-precision reality is
the point of \cref{rem:lec10-mult}).
\begin{enumerate}[leftmargin=1.6em,nosep]
  \item \emph{Phase estimation} (\cref{ch:lec04}) writes each eigenphase into a
    $d$-qubit register:
    \begin{equation}
      U_{\mathrm{QPE}}\ket{0^d}\ket b=\sum_j c_j\ket{\lambda_j}\ket{v_j}.
      \label{eq:lec10-qpe}
    \end{equation}
  \item A \emph{controlled rotation} on a fresh ancilla, driven by the eigenvalue
    register through coherent classical arithmetic $\lambda\mapsto\arcsin(C/\lambda)$,
    produces
    \begin{equation}
      U_{\mathrm{CR}}\ket0\ket{\lambda_j}
      =\Bigl(\sqrt{1-\tfrac{C^2}{\lambda_j^2}}\,\ket0+\tfrac{C}{\lambda_j}\ket1\Bigr)
        \ket{\lambda_j},
      \label{eq:lec10-cr}
    \end{equation}
    where $C\le\lambda_{\min}=1/\kappa$ is a normalization constant chosen so
    $\abs{C/\lambda_j}\le1$.
  \item \emph{Uncomputation} applies $U_{\mathrm{QPE}}\dg$, returning the eigenvalue
    register to $\ket{0^d}$ and leaving
    \begin{equation}
      U_{\mathrm{HHL}}\ket0\ket b
      =\sum_j c_j\Bigl(\sqrt{1-\tfrac{C^2}{\lambda_j^2}}\,\ket0
        +\tfrac{C}{\lambda_j}\ket1\Bigr)\ket{v_j}.
      \label{eq:lec10-hhl}
    \end{equation}
  \item \emph{Postselection} measures the ancilla; outcome $1$ yields the unnormalized
    vector $\ket{\tilde x\!\succ}=\sum_j (C/\lambda_j)c_j\ket{v_j}=C\,A^{-1}\ket b$, i.e.\
    the target solution up to the global factor $C$.
\end{enumerate}

\begin{figure}[t]
\centering
\begin{minipage}{0.92\textwidth}
\centering
\footnotesize
\setlength{\fboxsep}{8pt}
\fbox{\parbox{0.9\textwidth}{\centering
ancilla: \quad$\ket0\to$ \hfill $U_{\mathrm{CR}}$ \hfill $\to$ measure ($1=$ success)\\[2pt]
clock ($d$ qubits): \quad$\ket{0^d}\to U_{\mathrm{QPE}}\to\ \ket{\lambda_j}\ \to
   U_{\mathrm{QPE}}\dg\to\ket{0^d}$\\[2pt]
system: \quad$\ket b=\sum_j c_j\ket{v_j}$ carried through; output
   $\propto\sum_j (c_j/\lambda_j)\ket{v_j}$.}}
\end{minipage}
\caption{The HHL circuit: QPE writes the eigenphases into the clock register, a
controlled rotation attaches amplitude $C/\lambda_j$ to $\ket1$, uncomputation clears
the clock, and postselecting $\ket1$ leaves $C\,A^{-1}\ket b$ in the system register.}
\label{fig:lec10-hhl}
\end{figure}

\paragraph{Success probability and cost.}
With $\norm A=1$ so that $\lambda_{\min}=1/\kappa$, the largest admissible constant is
$C=1/\kappa$. The success probability is the squared norm of $\ket{\tilde x\!\succ}$,
\begin{equation}
  p(1)=\norm{\ket{\tilde x\!\succ}}^2=C^2\norm{A^{-1}\ket b}^2
      =\frac{\xi^2}{\kappa^2},\qquad \xi\defeq\norm{A^{-1}\ket b}\in[1,\kappa],
  \label{eq:lec10-hhl-prob}
\end{equation}
using $\norm{A^{-1}\ket b}\ge\norm{A^{-1}}^{-1}\cdots$; the trivial bound $\xi\ge1$
gives $p(1)=\Omega(\kappa^{-2})$ (worst case), while an $\Omega(1)$ overlap of $\ket b$
with the smallest-eigenvalue eigenvector gives $\xi=\Omega(\kappa)$ and $p(1)=\Omega(1)$
(best case). A single run is dominated by QPE, whose cost to resolve $\lambda_{\min}
=\Theta(1/\kappa)$ to the required precision is $\bigo(\kappa/\epsilon)$ (see
\cref{rem:lec10-mult}). Repeating for success costs $1/p(1)=\bigo(\kappa^2/\xi^2)$
rounds, so the total query complexity is
\begin{equation}
  \bigo\!\Bigl(\frac{\kappa}{\epsilon}\cdot\frac{\kappa^2}{\xi^2}\Bigr)
  =\bigo\!\Bigl(\frac{\kappa^3}{\xi^2\epsilon}\Bigr),
  \label{eq:lec10-hhl-cost}
\end{equation}
i.e.\ $\bigo(\kappa/\epsilon)$ in the best case and $\bigo(\kappa^3/\epsilon)$ in the
worst case. \emph{Amplitude amplification} (\cref{ch:lec03}) replaces the
$1/p(1)$-repetition overhead by its square root, boosting the worst case to
$\bigo(\kappa^2/\epsilon)$; \emph{variable-time} amplitude amplification (a refinement
that amplifies the many well-conditioned branches early and the few ill-conditioned
ones late) further reduces the $\kappa^2$ to $\kappa$, at the cost of a more delicate
analysis.

\begin{remark}[Multiplicative accuracy and the $1/\epsilon$ in QPE]
\label{rem:lec10-mult}
Because the estimated eigenvalue $\tilde\lambda_j=\lambda_j(1+e_j)$ enters through
$1/\tilde\lambda_j$, a \emph{multiplicative} error $\abs{e_j}\le\epsilon/4$ is needed to
guarantee \eqref{eq:lec10-qlsp}: writing $\ket{\tilde x\!\succ}-C\,x
=C\sum_j c_j(1/\tilde\lambda_j-1/\lambda_j)\ket{v_j}$ and bounding
$\abs{1/\tilde\lambda_j-1/\lambda_j}=\abs{e_j}/(\lambda_j\abs{1+e_j})
\le(\epsilon/2)/\lambda_j$ gives $\norm{\ket{\tilde x\!\succ}-C x}\le(C\epsilon/2)\norm x$,
whence $\norm{\ket{\tilde x}-\ket x}\le\epsilon$. Since $\lambda_{\min}=\Theta(1/\kappa)$,
resolving it to multiplicative precision $\epsilon$ means additive precision
$\epsilon/\kappa$, so $T_{\mathrm{QPE}}=\Omega(\kappa/\epsilon)$. This inverse-linear
dependence on $\epsilon$ is the chief weakness of HHL and the reason to prefer QSVT,
which needs only $\log(1/\epsilon)$.
\end{remark}

\section{QSVT solver}
\label{sec:lec10-3}

QSVT (\cref{ch:lec09}) inverts $A$ without ever estimating an eigenvalue: it applies a
single bounded polynomial directly to the singular values. From the SVD
$A=W\Sigma V\dg$,
\begin{equation}
  A^{-1}=V\Sigma^{-1}W\dg=g^{(\mathrm{SV})}(A\dg),\qquad g(x)=\tfrac1x\ \text{(odd)},
  \label{eq:lec10-qsvt-target}
\end{equation}
so we approximate the odd function $1/x$ on the spectral support $[1/\kappa,1]$ by a
bounded odd polynomial.

\begin{theorem}[QSVT linear-system solver]
\label{thm:lec10-qsvt}
Let $U_A\in\mathrm{BE}_{\alpha,m}(A)$ with condition number $\kappa$ (singular values in
$[{1/\kappa},1]$ after normalization), let $U_b$ prepare $\ket b$, and set
$\xi=\norm{A^{-1}\ket b}\in[1,\kappa]$. There is a QSVT solver producing $\ket{\tilde x}$
with $\norm{\ket{\tilde x}-\ket x}\le\epsilon$ and constant success probability, using
\begin{equation}
  \bigo\!\Bigl(\frac{\kappa^2}{\xi}\log\tfrac\kappa\epsilon\Bigr)
  \quad\text{queries to } U_A,U_A\dg,\qquad
  \bigo\!\Bigl(\frac{\kappa}{\xi}\Bigr)\ \text{queries to } U_b.
  \label{eq:lec10-qsvt-cost}
\end{equation}
In the worst case $\xi=1$ this is $\bigo(\kappa^2\log(\kappa/\epsilon))$; in the best
case $\xi=\Theta(\kappa)$ it is $\bigo(\kappa\log(\kappa/\epsilon))$.
\end{theorem}

\begin{proof}[Proof sketch]
By \cref{exa:lec09-inv} there is an odd polynomial $P$ of degree
$d=\bigo(\kappa\log(1/\epsilon'))$ with $\abs{P(x)-1/(\eta\kappa x)}\le\epsilon'$ on
$[1/\kappa,1]$ and $\abs P\le1$ on $[-1,1]$ (here $\eta=\Theta(1)$ enforces the global
bound). \cref{thm:lec09-qsvt} block-encodes $P^{(\mathrm{SV})}(A\dg)=VP(\Sigma)V\dg
\approx(\eta\kappa)^{-1}A^{-1}$ with $\bigo(d)$ queries. Applying it to $\ket{0^m}\ket b$
and reading the flagged block yields the unnormalized
$\ket{\tilde x\!\succ}\approx(\eta\kappa)^{-1}A^{-1}\ket b$, of norm
$\approx\xi/(\eta\kappa)$, so the flag succeeds with probability $\approx(\xi/\eta\kappa)^2$.
Amplitude amplification (\cref{ch:lec03}) boosts this to $\Omega(1)$ in
$\bigo(\eta\kappa/\xi)$ rounds, each calling $U_A,U_A\dg$ order $d$ times and $U_b,U_b\dg$
order once, giving \eqref{eq:lec10-qsvt-cost}. To meet the final tolerance $\epsilon$ in
\eqref{eq:lec10-qlsp} one sets the polynomial precision $\epsilon'=\Theta(\xi\epsilon/\kappa)$,
which only affects the logarithm; hence the $\log(\kappa/\epsilon)$.
\end{proof}

\paragraph{The inversion polynomial, made explicit.}
The workhorse construction (Childs--Kothari--Somma) starts from the smooth odd extension
\begin{equation}
  g_b(x)=\frac{1-(1-x^2)^{b}}{x},
  \label{eq:lec10-cks}
\end{equation}
an odd polynomial of degree $2b-1$ that agrees with $1/x$ up to
$(1-x^2)^b/\abs x\le\ee^{-bx^2}/\abs x$; choosing $b=\bigo(\kappa^2\log(1/\epsilon))$
makes the error $\epsilon$ on $[1/\kappa,1]$. Truncating the Chebyshev expansion of
$g_b$ and discarding its negligible high-degree tail reduces the effective degree to
$d=\bigo(\kappa\log(1/\epsilon))$; dividing by $\eta\kappa$ (with $\eta=4/3$) enforces
$\abs P\le1$. The resulting phase factors come from a QSP solver
(\cref{sec:lec09-5}); a convex-optimization construction yields the same degree with an
equioscillating error and is what production code uses. The essential point is the
\emph{linear-in-$\kappa$, logarithmic-in-$\epsilon$} degree---an exponential improvement
in $\epsilon$ over HHL, and the reason QSVT is the default solver whenever the block
encoding is available.

\begin{example}[{Poisson's equation: $\kappa\sim N^{2/d}$}]
\label{exa:lec10-poisson}
Consider the one-dimensional Poisson problem $-u''(r)=b(r)$ on $[0,1]$ with $u(0)=u(1)=0$,
discretized by central differences on a uniform grid $r_i=ih$, $h=1/(N+1)$. The
Laplacian becomes the tridiagonal Toeplitz matrix
\begin{equation}
  A=\frac1{h^2}\,\mathrm{tridiag}(-1,\,2,\,-1)\in\R^{N\times N}.
  \label{eq:lec10-lap}
\end{equation}
A Toeplitz tridiagonal matrix $\mathrm{tridiag}(b,a,b)$ has eigenvalues
$\lambda_k=a+2\abs b\cos\!\bigl(\tfrac{k\pi}{N+1}\bigr)$; with $a=2/h^2$, $b=-1/h^2$ the
extreme eigenvalues are $\lambda_{\max}\approx4/h^2$ and $\lambda_{\min}\approx\pi^2$,
so
\begin{equation}
  \kappa(A)=\frac{\lambda_{\max}}{\lambda_{\min}}\approx\frac{4}{h^2\pi^2}=\bigo(N^2).
  \label{eq:lec10-poisson-kappa}
\end{equation}
Thus even the QSVT solver costs $\widetilde{\bigo}(\kappa^2)=\widetilde{\bigo}(N^4)$
(worst case) or $\widetilde{\bigo}(\kappa)=\widetilde{\bigo}(N^2)$, with only
$\log(1/\epsilon)$ precision dependence---no advantage over classical sparse solvers for
a fixed $1$D grid. The situation changes in \emph{high dimension}: the $d$-dimensional
Laplacian is $A_d=\sum_{i=1}^{d}\id^{\otimes(i-1)}\otimes A\otimes\id^{\otimes(d-i)}$
with $\norm{A_d}\approx4d/h^2$ and $\norm{A_d^{-1}}^{-1}\approx d\pi^2$, so
$\kappa(A_d)\approx\kappa(A)=\bigo(N^2)$ is \emph{independent of $d$}. Since the total
grid has $N^d$ points, a solver whose cost is $\poly(d)\cdot\widetilde{\bigo}(\kappa)$
can outperform classical methods whose per-iteration cost $\bigo(N^d)$ is exponential
in $d$---provided the evolution oracle and state preparation are efficient and one wants
only a scalar functional of the solution. This dimensional scaling $\kappa\sim N^{2/d}$
(here $N_{\mathrm{tot}}=N^d$, so $\kappa\sim N_{\mathrm{tot}}^{2/d}$) is the standard
lens through which PDE-derived QLSP instances are judged.
\end{example}

\section{Adiabatic and quantum Zeno solvers}
\label{sec:lec10-4}

Both HHL and QSVT, even with amplitude amplification, retain a quadratic $\kappa$
dependence in the worst case, because the solution norm $\norm{A^{-1}\ket b}$ can be as
small as the subnormalization $1/\kappa$. The near-optimal and optimal solvers avoid
postselection entirely by recasting QLSP as an \emph{eigenpath-tracing} problem: build a
family of Hamiltonians whose zero-eigenstate deforms continuously from $\ket b$ to
$\ket x$, and follow it either by slow (adiabatic) evolution or by repeated projective
measurement (Zeno).

\paragraph{QLSP as a zero-eigenvalue eigenpath.}
Take $A$ Hermitian positive definite with spectrum in $[1/\kappa,1]$, and set
$Q_b=\id-\proj b$. Define
\begin{equation}
  H_0=\begin{pmatrix}0&Q_b\\ Q_b&0\end{pmatrix},\qquad
  H_1=\begin{pmatrix}0&A Q_b\\ Q_b A&0\end{pmatrix}.
  \label{eq:lec10-H0H1}
\end{equation}
The null space of $H_1$ is spanned by $\ket0\ket x$ (with $A\ket x\propto\ket b$, since
$Q_b A\ket x\propto Q_b\ket b=0$) and $\ket1\ket b$; the null space of $H_0$ is spanned
by $\ket0\ket b$ and $\ket1\ket b$. Interpolating
\begin{equation}
  H(f)=(1-f)H_0+f H_1,\qquad f=f(s)\colon[0,1]\to[0,1],\ f(0)=0,\ f(1)=1,
  \label{eq:lec10-Hf}
\end{equation}
with scheduling function $f$, the state $\ket0\ket b$ (a null vector of $H_0$) is
deformed into $\ket0\ket x$ (a null vector of $H_1$), which encodes the solution. The
key structural fact is that the relevant gap never closes and shrinks only near the
end.

\begin{lemma}[Gap of the interpolating Hamiltonian]
\label{lem:lec10-gap}
For $A\succeq\kappa^{-1}\id$ the spectral gap $\Delta(f)$ of $H(f)$ around its
zero eigenvalue satisfies
\begin{equation}
  \Delta(f)\ge 1-f+f/\kappa\ \ge\ 1/\kappa.
  \label{eq:lec10-gap}
\end{equation}
\end{lemma}

\begin{proof}
Write $H(f)=\bigl(\begin{smallmatrix}0&M Q_b\\ Q_b M&0\end{smallmatrix}\bigr)$ with
$M=(1-f)\id+fA$. An eigenvalue $\lambda$ of $H(f)$ satisfies
$0=\det\!\bigl(\lambda^2\id-Q_b M^2 Q_b\bigr)$ (block elimination, using that $\lambda\id$
commutes with the off-diagonal blocks). On the subspace orthogonal to $\ket b$,
$Q_b M^2 Q_b$ acts as $M^2$, so the smallest nonzero $\lambda^2$ obeys
\begin{equation}
  \Delta^2(f)=\!\!\inf_{\braket{b}{\varphi}=0,\,\norm\varphi=1}\!\!
    \bra\varphi M^2\ket\varphi
    \ \ge\ \inf_{\norm\varphi=1}\bra\varphi M^2\ket\varphi
    =\bigl((1-f)+f/\kappa\bigr)^2,
  \label{eq:lec10-gap-pf}
\end{equation}
since the eigenvalues of $M=(1-f)\id+fA$ lie in $[(1-f)+f/\kappa,\ (1-f)+f]$. Taking
square roots gives \eqref{eq:lec10-gap}; the minimum over $f\in[0,1]$ is $1/\kappa$ at
$f=1$.
\end{proof}

\paragraph{Adiabatic solvers: $\mathrm{AQC}(p)$ and $\mathrm{AQC}(\exp)$.}
Evolving $\tfrac{\ii}{T}\partial_s\ket{\psi_T(s)}=H(f(s))\ket{\psi_T(s)}$ from
$\ket{\psi_T(0)}=\ket0\ket b$, the adiabatic theorem guarantees that
$\ket{\psi_T(1)}$ lands in the null space of $H_1$---the solution---provided $T$ is
large relative to the gap. The naive bound scales the runtime as
$T=\bigo(\Delta_*^{-3}/\epsilon)=\bigo(\kappa^3/\epsilon)$, ruinous for small gaps. Two
ideas fix this.
\begin{itemize}[leftmargin=1.4em,nosep]
  \item \emph{Gap-adapted schedule} $\mathrm{AQC}(p)$: choose $\dot f(s)\propto
    \Delta^p(f(s))$ so the evolution slows exactly where the gap is small. For
    $1<p<2$ this yields $T=\bigo(\kappa/\epsilon)$---linear in $\kappa$, but still
    $1/\epsilon$.
  \item \emph{Boundary-cancelling schedule} $\mathrm{AQC}(\exp)$: choose
    \begin{equation}
      f(s)=c_e^{-1}\!\int_0^s\!\exp\!\Bigl(-\tfrac1{\tau(1-\tau)}\Bigr)\dd\tau,\qquad
      c_e=\!\int_0^1\!\exp\!\Bigl(-\tfrac1{\tau(1-\tau)}\Bigr)\dd\tau,
      \label{eq:lec10-aqcexp}
    \end{equation}
    which makes \emph{all} derivatives $H^{(k)}(0)=H^{(k)}(1)=0$ vanish at the
    endpoints. Integration by parts then suppresses the adiabatic error faster than any
    polynomial in $1/T$, giving runtime
    \begin{equation}
      T=\bigo\!\Bigl(\kappa\,\polylog\tfrac\kappa\epsilon\Bigr)
      \label{eq:lec10-aqcexp-cost}
    \end{equation}
    (near-optimal in $\kappa$, exponentially good in $\epsilon$). The continuous
    evolution \eqref{eq:lec10-Hf} must still be discretized, e.g.\ by a truncated Dyson
    series (\cref{ch:lec13}), contributing only polylogarithmic overhead.
\end{itemize}

\paragraph{The quantum Zeno / eigenstate-filtering solver.}
A digital alternative discretizes the schedule at $f_j=f(j/M)$ and, instead of
simulating a continuous evolution, \emph{projects} onto the successive zero-eigenspaces
by block-encoded spectral projectors $P_{f_j}$. Frequent projection freezes the state
on the eigenpath (the quantum Zeno effect: many measurements suppress transitions), and
only the last projection determines the final accuracy. Two ingredients make this
optimal.

First, the schedule
\begin{equation}
  f(s)=\frac{1-\kappa^{-s}}{1-\kappa^{-1}}
  \label{eq:lec10-zeno-schedule}
\end{equation}
equalizes the eigenpath length across steps: by perturbation theory
$\norm{\partial_f\ket{x(f)}}\le2/\Delta_*(f)$ with $\Delta_*(f)=1-f+f/\kappa$, so the
total path length is
\begin{equation}
  L(0,1)\le\int_0^1\frac{2}{1-f+f/\kappa}\,\dd f=\frac{2\log\kappa}{1-1/\kappa}
  =\bigo(\log\kappa),
  \label{eq:lec10-pathlen}
\end{equation}
and the schedule concentrates steps where the gap is small. With $M=\Theta(\log^2\kappa)$
steps the consecutive overlaps satisfy $\abs{\braket{x(f_j)}{x(f_{j-1})}}\ge
1-\tfrac12 L_*(f_{j-1},f_j)^2$, so all $M$ projections succeed with probability
$\Omega(1)$.

Second, each projector $P_{f_j}$ is realized by an \emph{eigenstate filter}: an even
polynomial approximating the rectangle that is $\approx1$ at eigenvalue $0$ and
$\approx0$ beyond the gap $\Delta_*(f_j)$, built by QSVT with degree
$\bigo(\Delta_*(f_j)^{-1}\log(1/\epsilon_j))$ (\cref{sec:lec09-3}, the Lin--Tong
filter). Summing the per-step costs along the schedule,
\begin{equation}
  \sum_{j=1}^{M-1}\frac{\log(1/\epsilon')}{1-f(s_j)+f(s_j)/\kappa}
  \ \lesssim\ M\log\tfrac1{\epsilon'}\!\int_0^1\!\kappa^{s}\,\dd s
  = M\log\tfrac1{\epsilon'}\,\frac{\kappa}{\log\kappa},
  \label{eq:lec10-zeno-sum}
\end{equation}
and taking $M=\Theta(\log\kappa)$ (with fixed-point amplitude amplification to keep each
step's success near $1$) plus a final projector of precision $\epsilon$ gives the
following optimal scaling.

\begin{theorem}[Quantum Zeno / eigenstate-filtering solver]
\label{thm:lec10-zeno}
Let $A$ have singular values in $[1/\kappa,1]$, accessed by
$U_A\in\mathrm{BE}_{\alpha,m}(A)$, and let $U_b$ prepare $\ket b$. Then a QZE-based
solver prepares $\ket{\tilde x}$ with $\abs{\braket{\tilde x}{x}}\ge1-\epsilon$ and
constant success probability, using
\begin{equation}
  \bigo\!\bigl(\kappa\log(1/\epsilon)\bigr)
  \label{eq:lec10-zeno-cost}
\end{equation}
queries to $U_A$ and $U_b$, up to $\log\log\kappa$ factors.
\end{theorem}

\begin{proof}[Proof idea]
Combine \cref{lem:lec10-gap}, the equalized schedule \eqref{eq:lec10-zeno-schedule}, the
path-length bound \eqref{eq:lec10-pathlen}, and the eigenstate-filter cost. The first
$M-1$ projectors cost $\bigo(\kappa(\log\log\kappa)^2)$ in total by
\eqref{eq:lec10-zeno-sum} with $M=\Theta(\log\kappa)$ and fixed-point amplification
holding each step's fidelity at $1-1/M$; the final projector costs
$\bigo(\kappa\log(1/\epsilon))$. Summing gives \eqref{eq:lec10-zeno-cost}. The matching
lower bound $\Omega(\kappa)$ is \cref{prop:lec10-lb}, so the scaling is optimal in
$\kappa$ and in $\epsilon$. A closely related \emph{discrete} adiabatic theorem
(Costa et al.) achieves the same $\bigo(\kappa\log(1/\epsilon))$ without the extra
$\log\log\kappa$.
\end{proof}

\section{Complexity landscape and preconditioning}
\label{sec:lec10-5}

\paragraph{The landscape.}
The methods above form a clean progression in the two parameters $\kappa$ and
$\epsilon$. The table below collects the query complexities to the oracle
for $A$ (worst case, before/after amplitude amplification as noted).

\begin{center}
\renewcommand{\arraystretch}{1.3}
\begin{tabular}{@{}lll@{}}
\hline
\textbf{Method} & \textbf{Queries to $A$} & \textbf{Precision dep.}\\
\hline
HHL & $\bigo(\kappa^2/\epsilon)$ (with AA) & $1/\epsilon$\\
LCU (Fourier/Chebyshev) & $\bigo(\kappa^2\polylog(\kappa/\epsilon))$ & $\polylog$\\
QSVT & $\bigo(\kappa^2\log(\kappa/\epsilon))$ & $\log(1/\epsilon)$\\
$\mathrm{AQC}(p)$ & $\bigo((\kappa/\epsilon)\log(\kappa/\epsilon))$ & $1/\epsilon$\\
$\mathrm{AQC}(\exp)$ & $\bigo(\kappa\,\polylog(\kappa/\epsilon))$ & $\polylog$\\
QZE $+$ eigenfilter & $\bigo(\kappa\log(1/\epsilon))$ & $\log(1/\epsilon)$\\
discrete AQC (optimal) & $\bigo(\kappa\log(1/\epsilon))$ & $\log(1/\epsilon)$\\
\hline
\end{tabular}
\end{center}

\noindent\emph{Table:} QLSP query complexity in the worst case
$\xi=\norm{A^{-1}\ket b}=1$. Best case ($\xi=\Theta(\kappa)$) removes one factor of
$\kappa$ from the postselection methods. The optimal $\bigo(\kappa\log(1/\epsilon))$
matches the lower bound $\Omega(\kappa)$ of \cref{prop:lec10-lb}.

Two lessons. First, the frontier moved from HHL's $\kappa^2/\epsilon$ to the optimal
$\kappa\log(1/\epsilon)$ by two independent improvements---better \emph{$\epsilon$
dependence} (QSVT's polynomial replacing QPE's arithmetic) and better \emph{$\kappa$
dependence} (eigenpath methods replacing postselection). Second, even the optimal solver
is $\Omega(\kappa)$: unlike classical conjugate gradient's $\bigo(\sqrt\kappa)$ for
positive-definite systems, generic quantum solvers do not achieve $\sqrt\kappa$, and no
amount of algorithmic cleverness removes the linear $\kappa$. Hence, whenever $\kappa$ is
large, the decisive move is to \emph{reduce $\kappa$ itself}---preconditioning.

\paragraph{Why $\kappa$, not $N$, is the enemy.}
The comparison with classical iterative solvers is instructive. For a $d$-sparse
Hermitian positive-definite $A$, one matrix--vector product costs $\bigo(dN)$ flops, and
steepest descent needs $\bigo(\kappa\log(1/\epsilon))$ iterations, conjugate gradient
$\bigo(\sqrt\kappa\log(1/\epsilon))$; total classical cost
$\bigo(dN\sqrt\kappa\log(1/\epsilon))$. The quantum solver's cost is
$\widetilde{\bigo}(\kappa)$ \emph{queries}, each of which may be a $\poly(n)$-gate
circuit when $U_A$ is efficiently implementable---so the potential exponential advantage
is in replacing the $\bigo(N)$ per-step cost by $\poly(\log N)$, \emph{not} in the
$\kappa$ dependence, which is if anything worse than CG. The advantage is real only when
(i) $U_A$ and $U_b$ are cheap, (ii) one wants a scalar functional of $\ket x$ rather than
all of $x$, and (iii) $\kappa$ is controlled. The Poisson example
(\cref{exa:lec10-poisson}) with $\kappa\sim N^{2/d}$ shows all three tensions at once.

\paragraph{Preconditioning: classical and quantum.}
Left-preconditioning replaces $Ax=b$ by $MAx=Mb$; it helps when $\kappa(MA)\ll\kappa(A)$.
Classical choices---diagonal (Jacobi), incomplete factorizations, sparse approximate
inverse (SPAI)---have quantum analogues (SPAI, circulant preconditioners), but a genuine
obstacle appears: on a quantum computer every block-encoded matrix has spectral norm
$\le1$, so multiplying by $M$ can only \emph{shrink} the state, and recovering the norm
costs $\bigo(\kappa)$ amplitude-amplification rounds. Applying two preconditioners
therefore incurs a \emph{subnormalization} $\kappa(M)\kappa(M_r)\propto\kappa(A)$, which
can eat the very benefit sought. Concretely, the complexity of a quantum solve is
governed not by $\kappa(S)$ alone but by $\tilde\gamma(U_S)\,\kappa(S)$, the product of
the block-encoding subnormalization and the condition number. A clean pipeline where
this works: if $A+B$ is the target with $A$ easily invertible and of large norm and
$B,A^{-1},(A+B)^{-1}$ moderate, then $\kappa(\id+A^{-1}B)=\bigo(1)$, and one composes
block encodings along
\begin{equation}
  A^{-1}\ \to\ A^{-1}B\ \to\ \id+A^{-1}B\ \to\ (\id+A^{-1}B)^{-1}
  \ \to\ (\id+A^{-1}B)^{-1}A^{-1}=(A+B)^{-1},
  \label{eq:lec10-precond-pipeline}
\end{equation}
requiring only matrix addition, multiplication (\cref{ch:lec07}), and fast inversion of
the well-conditioned factor.

\paragraph{Optimal multilevel (BPX) preconditioning.}
For elliptic PDEs the sharpest tool is a \emph{multilevel} preconditioner. Discretizing
$-\operatorname{div}(A\nabla u)=f$ by finite elements on a hierarchy of grids
$G_1\subset\cdots\subset G_L$ with $h_\ell=2^{-\ell}$ gives a stiffness matrix $S$ with
$\kappa(S)=\bigo(h^{-2})=\bigo(N^{2/d})$. The Bramble--Pasciak--Xu (BPX) preconditioner
$P$ is a hierarchical decomposition across levels; under the standard Jackson (direct)
and Bernstein (inverse) estimates,
\begin{equation}
  \inf_{v_\ell\in V_\ell}\norm{u-v_\ell}_{L^2}\le 2^{-\ell}C_1\norm{u}_{H^1},\qquad
  \norm{v_\ell}_{H^1}\le 2^{\ell}C_2\norm{v_\ell}_{L^2},
  \label{eq:lec10-jackson-bernstein}
\end{equation}
the symmetrically preconditioned matrix satisfies $\kappa(P^{1/2}SP^{1/2})=\bigo(1)$,
\emph{independent of $L$}. Because the preconditioned system $PS$ itself need not be well
conditioned as a product, one \emph{splits} $P=FF^{\mathsf T}$ over a multilevel
generating system $\mathcal F=\{2^{-\ell(2-d)/2}\lambda_j^{(\ell)}\}$ and solves the
frame system $F^{\mathsf T}SF\,y=F^{\mathsf T}r$, $c=Fy$. On Cartesian grids the
resulting block encoding of $F^{\mathsf T}SF$ has subnormalization $\bigo(dL)$ and
$\kappa=\bigo(1)$, so the quantum finite-element solve attains query complexity
near-optimal in $\kappa$ and $\epsilon$---turning a $\kappa\sim N^{2/d}$ instance into an
$\bigo(1)$-conditioned one and making the QLSP speedup accessible for elliptic PDEs. This
is the state of the art for reconciling the $\Omega(\kappa)$ lower bound with practical
PDE problems.

\exercisesection

\begin{exercise}\label{ex:lec10-1}
\emph{Condition number and normalization.} Let $A\in\C^{N\times N}$ be invertible with
$\norm A\le1$. Show that the singular values of $A$ lie in $[1/\kappa,1]$ where
$\kappa=\norm A\norm{A^{-1}}$, and that the Hermitian dilation
$\widetilde A=\bigl(\begin{smallmatrix}0&A\\ A\dg&0\end{smallmatrix}\bigr)$ has
eigenvalues $\pm\sigma_i$ and the same condition number. Verify that
$\widetilde A(0,x)^{\mathsf T}=(b,0)^{\mathsf T}$ is equivalent to $Ax=b$.
\end{exercise}

\begin{exercise}\label{ex:lec10-2}
\emph{HHL success probability.} From \eqref{eq:lec10-hhl} with $C=1/\kappa$, show the
postselection success probability is $p(1)=\norm{A^{-1}\ket b}^2/\kappa^2=\xi^2/\kappa^2$.
Deduce $p(1)=\Omega(\kappa^{-2})$ using only $\xi\ge1$, and $p(1)=\Omega(1)$ if $\ket b$
has $\Omega(1)$ overlap with the eigenvector of the smallest eigenvalue. Explain how
amplitude amplification turns the $\kappa^2$ repetition overhead into $\kappa$.
\end{exercise}

\begin{exercise}\label{ex:lec10-3}
\emph{Multiplicative precision in HHL.} Suppose the eigenvalue estimates satisfy
$\tilde\lambda_j=\lambda_j(1+e_j)$ with $\abs{e_j}\le\epsilon/4$. Following
\cref{rem:lec10-mult}, show $\norm{\ket{\tilde x}-\ket x}\le\epsilon$, and conclude that
because $\lambda_{\min}=\Theta(1/\kappa)$ the QPE runtime must be $\Omega(\kappa/\epsilon)$.
Contrast with the $\log(1/\epsilon)$ dependence of the QSVT solver.
\end{exercise}

\begin{exercise}\label{ex:lec10-4}
\emph{The QSVT inversion polynomial.} For $g_b(x)=(1-(1-x^2)^b)/x$ of
\eqref{eq:lec10-cks}, show $g_b$ is an odd polynomial of degree $2b-1$ and that
$\abs{g_b(x)-1/x}\le\ee^{-bx^2}/\abs x$ for $x\neq0$. Deduce that
$b=\bigo(\kappa^2\log(1/\epsilon))$ gives error $\epsilon$ on $[1/\kappa,1]$, and explain
why Chebyshev truncation reduces the effective QSVT degree to
$\bigo(\kappa\log(1/\epsilon))$. Why must one divide by $\eta\kappa$ with $\eta>1$?
\end{exercise}

\begin{exercise}\label{ex:lec10-5}
\emph{Gap of the interpolating Hamiltonian.} Prove \cref{lem:lec10-gap}: for
$A\succeq\kappa^{-1}\id$, the gap of $H(f)=(1-f)H_0+fH_1$ around its zero eigenvalue is
$\Delta(f)\ge1-f+f/\kappa$. (Use the block-elimination identity
$\det(\lambda^2\id-Q_bM^2Q_b)=0$ with $M=(1-f)\id+fA$, and restrict to the subspace
orthogonal to $\ket b$.) Where is the gap smallest, and what value does it take?
\end{exercise}

\begin{exercise}\label{ex:lec10-6}
\emph{Eigenpath length and the Zeno schedule.} For the schedule
$f(s)=(1-\kappa^{-s})/(1-\kappa^{-1})$ of \eqref{eq:lec10-zeno-schedule} and the bound
$\norm{\partial_f\ket{x(f)}}\le2/(1-f+f/\kappa)$, show that the total eigenpath length is
$L(0,1)\le 2\log\kappa/(1-1/\kappa)=\bigo(\log\kappa)$. Explain why this schedule
allocates roughly equal path length per step, and why that yields the
$\bigo(\kappa\log(1/\epsilon))$ query complexity of \cref{thm:lec10-zeno}.
\end{exercise}

\begin{exercise}\label{ex:lec10-7}
\emph{Poisson condition number.} Using the eigenvalues
$\lambda_k=a+2\abs b\cos(k\pi/(N+1))$ of a Toeplitz tridiagonal matrix, derive
$\kappa(A)\approx4/(h^2\pi^2)=\bigo(N^2)$ for the $1$D Laplacian
\eqref{eq:lec10-lap}. Show that the $d$-dimensional Laplacian
$A_d=\sum_i\id^{\otimes(i-1)}\otimes A\otimes\id^{\otimes(d-i)}$ has
$\kappa(A_d)\approx\kappa(A)$ independent of $d$, and explain why this makes
high-dimensional Poisson a candidate for quantum advantage while the $1$D case is not.
\end{exercise}

\begin{exercise}\label{ex:lec10-8}
\emph{The $(A+B)$ preconditioner.} Suppose $A$ is easily invertible with large
$\norm A$, and $B,\norm{A^{-1}},\norm{(A+B)^{-1}}$ are moderate. Show
$\kappa(\id+A^{-1}B)=\bigo(1)$ and verify the identity
$(\id+A^{-1}B)^{-1}A^{-1}=(A+B)^{-1}$. Using the block-encoding arithmetic of
\cref{ch:lec07}, describe the pipeline \eqref{eq:lec10-precond-pipeline} and identify at
which step the well-conditioned inversion (cheap) versus the ill-conditioned one
(expensive) occurs. Why does naive left-preconditioning $M A$ incur a subnormalization
$\propto\kappa(M)$?
\end{exercise}

\readingsection
\begin{itemize}[leftmargin=1.4em]
  \item Lin Lin, \emph{live notes}, Ch.~19 (QLSP problem setup and output model; the HHL
        algorithm via QPE and controlled rotation, its success probability and
        multiplicative-accuracy requirement; the QSVT solver and
        \cref{thm:lec10-qsvt}; Poisson's equation and fermionic Green's functions; the
        quantum adiabatic near-optimal solver $\mathrm{AQC}(\exp)$; and the quantum Zeno
        / eigenstate-filtering optimal solver, \cref{thm:lec10-zeno}).
  \item Dong An, \emph{Quantum Numerical Linear Algebra} tutorial talk (the QLSP
        definition and parameters $\epsilon,\kappa$; QSVT for $f(x)=1/(\kappa x)$ with
        degree $\bigo(\kappa\log(1/\epsilon))$; HHL as QPE $+$ controlled rotation with
        cost $\bigo(\kappa^2/\epsilon)$; the LCU Fourier/Chebyshev approaches; the
        $\mathrm{AQC}(p)$/$\mathrm{AQC}(\exp)$ schedules and their
        $\bigo(\kappa/\epsilon)$ and $\bigo(\kappa\,\polylog)$ costs; the complexities
        comparison table; and preconditioning, including SPAI and the $(A+B)$ pipeline).
  \item A.~W.~Harrow, A.~Hassidim, and S.~Lloyd, \emph{Quantum algorithm for linear
        systems of equations}, Phys.\ Rev.\ Lett.\ \textbf{103}, 150502 (2009)
        (arXiv:0811.3171), the original HHL algorithm.
  \item A.~M.~Childs, R.~Kothari, and R.~D.~Somma, \emph{Quantum algorithm for systems
        of linear equations with exponentially improved dependence on precision}, SIAM
        J.\ Comput.\ \textbf{46}, 1920--1950 (2017) (arXiv:1511.02306), for the LCU
        solver and the explicit $1/x$ polynomial; D.~An and L.~Lin, \emph{Quantum linear
        system solver based on time-optimal adiabatic quantum computing and quantum
        approximate optimization algorithm}, ACM Trans.\ Quantum Comput.\ \textbf{3}, 5
        (2022) (arXiv:1909.05500), for $\mathrm{AQC}(p)$ and $\mathrm{AQC}(\exp)$;
        L.~Lin and Y.~Tong, \emph{Optimal polynomial based quantum eigenstate filtering
        with application to solving quantum linear systems}, Quantum \textbf{4}, 361
        (2020) (arXiv:1910.14596), for the QZE/eigenstate-filtering optimal solver; and
        P.~C.~S.~Costa et al., \emph{Optimal scaling quantum linear-systems solver via
        discrete adiabatic theorem}, PRX Quantum \textbf{3}, 040303 (2022)
        (arXiv:2111.08152).
  \item Quantum preconditioning notes (SuzhouNotes, ``Quantum preconditioning'') for the
        block-encoding normalization/subnormalization accounting, the QLSP definition in
        terms of $\tilde\gamma\cdot\kappa$, and the optimal multilevel (BPX)
        preconditioner following J.~H.~Bramble, J.~E.~Pasciak, and J.~Xu,
        \emph{Parallel multilevel preconditioners}, Math.\ Comp.\ \textbf{55}, 1--22
        (1990).
\end{itemize}
